% ---
% id: struve-2023-reg-7-1
% kind: solution
% lang: ru
% ---
Марс с Земли — не можем. Марс имел бы вид тонкого серпа, если бы мы смотрели
на него с теневой стороны. Но Марс — внешняя планета, то есть находится
дальше от Солнца, чем Земля. Поэтому мы всегда наблюдаем ту сторону Марса,
которая обращена к Солнцу, то есть освещённую.

Землю с Марса — можем. Принципиально наблюдение Земли с Марса не отличается
от наблюдения Венеры или Меркурия (внутренних планет) с Земли. Однако
максимальный видимый угловой размер Земли с Марса примерно равен
максимальному видимому угловому размеру Венеры с Земли (радиусы сопоставимые,
расстояния тоже), так что потребуется телескоп (или бинокль).

\medskip\noindent\textbf{Критерии оценивания}
\begin{center}
\begin{tabular}{@{}llr@{}}
  а1 & Верный ответ: Марс с Земли — не можем & 1 \\
  а2 & Корректное (по смыслу) обоснование ответа & 3 \\
  б1 & Верный ответ: Землю с Марса — можем & 1 \\
  б2 & Корректное (по смыслу) обоснование ответа & 1 \\
  б3 & Упоминание про необходимость телескопа или бинокля & 2 \\
  \multicolumn{2}{@{}l}{\textbf{Всего}} & \textbf{8} \\
\end{tabular}
\end{center}
